Large language models are increasingly deployed in applications requiring up-to-date knowledge, yet their parametric knowledge is frozen at training time. We investigate \emph{which types} of post-cutoff text resist fine-tuning and why, moving beyond the established observation that temporal degradation exists to characterize the learning dynamics that underlie it. Using \mistral (September 2023 cutoff) with \lora fine-tuning on domain-stratified text spanning 2022--2025 and API-based knowledge probing of \gptthree, we find a clear dichotomy: \textbf{novel factual content is stuck} while \textbf{textual patterns are flexible}. News text exhibits a statistically significant linear perplexity increase of 0.065 PPL/month with temporal distance ($R^2{=}0.935$, $p{=}0.033$), and fine-tuning closes only 40\% of this gap. In contrast, structured factual text achieves 75\% loss reduction regardless of whether content is pre- or post-cutoff, with the temporal gap shrinking by 80\% after training. API probing confirms these findings: \gptthree scores only 13.3\% on post-cutoff factual questions, with correct answers limited to events partially known before its cutoff. Our results show that fine-tuning teaches models to \emph{sound like} recent text without learning the facts within it, with direct implications for deployment: use retrieval for factual recency and fine-tuning for domain adaptation.
